%% Question bank chapter file - generated by bank_tools.py build (do not edit while the bank is still being built:
%% it is regenerated from the per-test data in Question Bank\data). Every question is one bankq block:
%%   \begin{bankq}[id=..., source={...}, paper=.., qno=.., type=..., subtopic={...}, answer={...}, answer_src=key|solved,
%%                 working={...}, difficulty=1-10, diff_parts={K S I C F}, diff_note={...},
%%                 group=<passage id or empty>, also_in={duplicate copies in other tests}]
%% difficulty = min(10, K+S+I+C+F-1): K concepts 1-3, S steps 1-4, I insight 0-2, C computation 0-2, F format 0-1
%%   ...question LaTeX, options, figures (\bfig / \bfigi -> img/)...
%%   \end{bankq}
%% Passages shared by several questions are bankpassage blocks; their questions carry group=<passage id>.
\documentclass[10pt]{article}
\usepackage[a4paper,margin=1.8cm]{geometry}
\usepackage[T1]{fontenc}
\usepackage{amsmath,amssymb,mathtools}
\usepackage{graphicx}
\usepackage[export]{adjustbox}
\usepackage{xcolor,needspace,keyval}
\usepackage{enumitem}
\usepackage{booktabs,array,tabularx}
\usepackage[hidelinks]{hyperref}
\setlength{\parindent}{0pt}
\setlength{\parskip}{3pt}
\makeatletter
\define@key{bq}{id}{\def\bq@id{#1}}\define@key{bq}{source}{\def\bq@src{#1}}\define@key{bq}{paper}{\def\bq@paper{#1}}
\define@key{bq}{qno}{\def\bq@qno{#1}}\define@key{bq}{type}{\def\bq@type{#1}}\define@key{bq}{subtopic}{\def\bq@sub{#1}}
\define@key{bq}{answer}{\def\bq@ans{#1}}\define@key{bq}{answer_src}{\def\bq@asrc{#1}}\define@key{bq}{working}{\def\bq@work{#1}}
\define@key{bq}{group}{\def\bq@group{#1}}\define@key{bq}{also_in}{\def\bq@also{#1}}
\define@key{bq}{difficulty}{\def\bq@diff{#1}}\define@key{bq}{diff_parts}{\def\bq@dparts{#1}}\define@key{bq}{diff_note}{\def\bq@dnote{#1}}
\newenvironment{bankq}[1][]{\def\bq@work{}\def\bq@also{}\def\bq@group{}\def\bq@diff{?}\def\bq@dparts{}\def\bq@dnote{}\setkeys{bq}{#1}%
  \par\needspace{4\baselineskip}\medskip\noindent{\small\textbf{\bq@id}\hfill\color{gray}[\bq@src, Paper \bq@paper, Q\bq@qno]
  \bq@type\ $\cdot$ \bq@sub\ $\cdot$ \textbf{D\bq@diff}}\par\nopagebreak\relax}
  {\par{\small\color{teal}Answer: \textbf{\bq@ans}\ifx\bq@work\@empty\else\ (solved: \bq@work)\fi
  \ifx\bq@also\@empty\else\par\color{gray}Also in: \bq@also\fi}\par}
\define@key{bp}{id}{\def\bp@id{#1}}\define@key{bp}{source}{\def\bp@src{#1}}\define@key{bp}{questions}{\def\bp@qs{#1}}
\newenvironment{bankpassage}[1][]{\setkeys{bp}{#1}\par\needspace{4\baselineskip}\bigskip\noindent
  {\small\textbf{Passage \bp@id}\hfill\color{gray}[\bp@src; questions \bp@qs]}\par\nopagebreak\relax\itshape}{\par}
\makeatother
\newcommand{\bfig}[2]{\par\smallskip{\centering\includegraphics[width=#1pt,max width=0.62\linewidth,max height=0.32\textheight]{img/#2}\par}\smallskip}
\newcommand{\bfigi}[2]{\includegraphics[width=#1pt,max width=\linewidth,max height=0.32\textheight,valign=c]{img/#2}}
\newcommand{\bmissing}[1]{\textit{[figure missing in the source conversion: \detokenize{#1}]}}
\newenvironment{qoptions}{\begin{itemize}[leftmargin=1.8em,itemsep=1pt,topsep=2pt,label={}]}{\end{itemize}}
\newcommand{\opt}[2]{\item[(#1)] #2}
\providecommand{\fwstart}{}\providecommand{\fwend}{}\providecommand{\qfullsplit}{}
\newcolumntype{L}[1]{>{\raggedright\arraybackslash}p{#1}}
\newcolumntype{C}[1]{>{\centering\arraybackslash}p{#1}}
\newcolumntype{R}[1]{>{\raggedleft\arraybackslash}p{#1}}
\providecommand{\sectionhdr}[1]{}\providecommand{\subsectionhdr}[1]{}


\begin{document}
{\Large\bfseries p-block --- Question Bank}\par\smallskip
{\small Chemistry chapter C36. Questions from all converted test papers, grouped by source series. Answers: printed key where it matches the paper, otherwise solved (answer\_src=solved, with working). Exact repeats across papers appear once, with also\_in listing the other copies.}\par\medskip\hrule

\section*{Score-I papers}
\begin{bankq}[id=SI12-P1-CHEM-S3-Q08, source={Score-I - Test 12 (09 Mar 2022)}, paper=1, qno=8, type=numerical, subtopic={Group 14}, answer={6}, answer_src=key, difficulty=3, diff_parts={K1 S2 I0 C0 F1}, diff_note={concepts: thermal decomposition reactions | steps: PbCl4 and Mg(NO3)2; HgC2O4 and Ag2CO3; Pb(CN)4 and Cu(CN)2; Al(OH)3 is not redox}]
Find the total number of compounds which on heating undergo redox reactions.\\
PbCl$_4$, Mg(NO$_3$)$_2$, HgC$_2$O$_4$, Ag$_2$CO$_3$, Pb(CN)$_4$, Al(OH)$_3$, Cu(CN)$_2$
\end{bankq}
\begin{bankq}[id=SI12-P1-CHEM-S3-Q10, source={Score-I - Test 12 (09 Mar 2022)}, paper=1, qno=10, type=numerical, subtopic={Group 18}, answer={3}, answer_src=key, difficulty=3, diff_parts={K1 S2 I0 C0 F1}, diff_note={concepts: VSEPR shapes of XeF4 and XeF5- | steps: P is the XeF5- salt; angles 90 and 72; atoms in plane 6 then (90 - 72)/6}]
Consider the following interaction\\
NaF + XeF$_4\to$ P (ionic compound)\\
Find the value of $\left[\dfrac{X-Y}{Z}\right]$\\
where\\
$X$ = Value of adjacent bond angle in XeF$_4$\\
$Y$ = Value of adjacent bond angle in anionic part of compound `P'\\
$Z$ = Maximum number of atoms that can lie in a plane in anionic part of compound `P'
\end{bankq}

\section*{TOAS papers}
\begin{bankq}[id=TOAS02-P1-CHEM-S2-Q12, source={TOAS - Test 02 (13 Feb 2022)}, paper=1, qno=12, type=multi, subtopic={Group 15}, answer={A,B,D}, answer_src=solved, working={AB$_3$ = NH$_3$ (unsymmetric cleavage of B$_2$H$_6$, dipole larger than NCl$_3$); AC$_3$ = NCl$_3$; C$_2$ = Cl$_2$, BC = HCl (A true); N-H < N-Cl (B true); NCl$_3$ has no back bonding (C false); N forms neutral oxides NO, N$_2$O (D true)}, difficulty=5, diff_parts={K1 S3 I1 C0 F1}, diff_note={concepts: identification of N compounds from properties | steps: identify NH3 from the diborane cleavage; identify NCl3 and its hydrolysis; compare bond lengths; back bonding; oxides of N | I: identification from clues}]
The compound $\mathrm{AB_3}$ has more dipole moment than $\mathrm{AC_3}$. The compound $\mathrm{AC_3}$ on action of superheated steam under high pressure yield an acidic oxide via $\mathrm{S_N1}$ mechanism. $\mathrm{AB_3}$ has tendency to unsymmetrically cleave bridge bond of $\mathrm{B_2H_6}$. Which is/are correct in view of above information.
\begin{qoptions}
\opt{A}{The compound ``BC'' can be produced by hydrolysis of $\mathrm{C_2}$.}
\opt{B}{The compound $\mathrm{BA_3}$ has less B--A bond length compared to C--A bond length in $\mathrm{CA_3}$.}
\opt{C}{$\mathrm{AC_3}$ show back bonding.}
\opt{D}{A can form neutral oxide}
\end{qoptions}
\end{bankq}
\begin{bankq}[id=TOAS03-P1-CHEM-S2-Q11, source={TOAS - Test 03 (27 Feb 2022)}, paper=1, qno=11, type=multi, subtopic={Group 17}, answer={A,C,D}, answer_src=key, difficulty=3, diff_parts={K1 S2 I0 C0 F1}, diff_note={concepts: structure and properties of chlorine oxides | steps: bond length in the two oxides; bond angles; colour of ClO2; solubility}]
Between $\mathrm{ClO_2}$ and $\mathrm{Cl_2O}$
\begin{qoptions}
\opt{A}{The Cl - O bond length would be larger in $\mathrm{Cl_2O}$}
\opt{B}{The bond angle would be larger in $\mathrm{Cl_2O}$}
\opt{C}{$\mathrm{ClO_2}$ is likely to be a coloured compound}
\opt{D}{$\mathrm{Cl_2O}$ is likely to be soluble in water}
\end{qoptions}
\end{bankq}
\begin{bankq}[id=TOAS03-P1-CHEM-S3-Q06, source={TOAS - Test 03 (27 Feb 2022)}, paper=1, qno=6, type=numerical, subtopic={Group 17}, answer={8}, answer_src=key, difficulty=2, diff_parts={K1 S1 I0 C0 F1}, diff_note={concepts: hydrolysis of interhalogens | steps: n = 5 for oxidation state +5; hydrolysis equation gives P = 3}]
$\mathrm{XY_n}$ is an interhalogen compound which produces oxyacid of $+5$ oxidation state of central atom on complete hydrolysis. If `P' moles $\mathrm{H_2O}$ is required for the complete hydrolysis of one mole $\mathrm{XY_n}$ then the value of $(n+P)$ is :
\end{bankq}
\begin{bankq}[id=TOAS03-P2-CHEM-S1-Q06, source={TOAS - Test 03 (27 Feb 2022)}, paper=2, qno=6, type=multi, subtopic={Group 13}, answer={A,C}, answer_src=key, difficulty=4, diff_parts={K1 S2 I1 C0 F1}, diff_note={concepts: boron hydrides and borazine reactions | steps: identify the products; aromatic character and boiling points; hybridisation of boron; back bonding | I: product identification}]
Consider the following reaction and choose correct statements
\bfig{265.2}{TOAS03_p2_chem_s1_q06_1.png}

\begin{qoptions}
\opt{A}{Product (B) contain aromatic character.}
\opt{B}{Boiling point of product (C) is less than Helium.}
\opt{C}{All Boron in product (A) are $sp^3$ hybridised.}
\opt{D}{In product (B), $2p_\pi$--$3p\pi$ type back bond present.}
\end{qoptions}
\end{bankq}
\begin{bankq}[id=TOAS03-P2-CHEM-S2-Q06, source={TOAS - Test 03 (27 Feb 2022)}, paper=2, qno=6, type=numerical, subtopic={Group 15}, answer={5}, answer_src=solved, working={PCl$_3\to$H$_3$PO$_3$ (dibasic, 2); SF$_4\to$H$_2$SO$_3$ (2); BrF$_5\to$HBrO$_3$ (1); sum 5}, difficulty=2, diff_parts={K1 S1 I0 C0 F1}, diff_note={concepts: hydrolysis of non-metal halides | steps: oxyacid formed from each compound; basicity of each and the sum}]
Write sum of basicity of oxy acid formed by hydrolysis of :-
\bfig{134.5}{TOAS03_p2_chem_s2_q06_1.png}
\bfig{134.2}{TOAS03_p2_chem_s2_q06_2.png}
\bfig{134.4}{TOAS03_p2_chem_s2_q06_3.png}
\end{bankq}
\begin{bankq}[id=TOAS04-P1-CHEM-S1-Q01, source={TOAS - Test 04 (15 Mar 2022)}, paper=1, qno=1, type=single, subtopic={Group 13}, answer={D}, answer_src=key, difficulty=1, diff_parts={K1 S1 I0 C0 F0}, diff_note={concepts: cleavage of diborane | steps: symmetrical cleavage with NMe3, THF and CO; unsymmetrical cleavage with NH3}]
In which of the following reaction bond cleavage of $\mathrm{B_2H_6}$ is different from the other reactions.
\begin{qoptions}
\opt{A}{$\mathrm{B_2H_6+2NMe_3\longrightarrow}$}
\opt{B}{$\mathrm{B_2H_6+2THF\longrightarrow}$}
\opt{C}{$\mathrm{B_2H_6+2CO\longrightarrow}$}
\opt{D}{$\mathrm{B_2H_6+2NH_3\longrightarrow}$}
\end{qoptions}
\end{bankq}
\begin{bankq}[id=TOAS05-P1-CHEM-S2-Q11, source={TOAS - Test 05 (03 Apr 2022)}, paper=1, qno=11, type=single, subtopic={Group 13}, answer={D}, answer_src=key, difficulty=2, diff_parts={K1 S2 I0 C0 F0}, diff_note={concepts: properties of boron trihalides | steps: acidic strength; bond length; hybridisation; match for X = F}]
{\small\setlength{\tabcolsep}{3pt}\renewcommand{\arraystretch}{1.25}\begin{center}\begin{tabular}{|L{0.8cm}|L{3.0cm}|L{0.65cm}|L{1.6cm}|}
\hline
\multicolumn{2}{|c|}{\parbox{3.3cm}{\centering List-I\\Dimerization/ polymerization by}} & \multicolumn{2}{c|}{List-II}\\
\hline
(I) & $3c-2e^-$ & (P) & $\mathrm{BeX_2}$\\
\hline
(II) & Co-ordinate bond & (Q) & $\mathrm{AlX_3}$\\
\hline
(III) & Odd $e^-$ & (R) & $\mathrm{CX_3}$\\
\hline
(IV) & No dimerization & (S) & $\mathrm{BX_3}$\\
\hline
\multicolumn{2}{|c|}{} & (T) & $\mathrm{IX_3}$\\
\hline
\multicolumn{2}{|c|}{} & (U) & $\mathrm{ClO_3}$\\
\hline
\end{tabular}\end{center}}
Select correct match when $X=\mathrm F$.
\begin{qoptions}
\opt{A}{II $\to$ R}
\opt{B}{I $\to$ Q}
\opt{C}{III $\to$ U}
\opt{D}{IV $\to$ S}
\end{qoptions}
\end{bankq}
\begin{bankq}[id=TOAS05-P1-CHEM-S2-Q12, source={TOAS - Test 05 (03 Apr 2022)}, paper=1, qno=12, type=single, subtopic={Group 13}, answer={C}, answer_src=key, difficulty=2, diff_parts={K1 S2 I0 C0 F0}, diff_note={concepts: properties of boron trihalides | steps: acidic strength; bond length; hybridisation; match for X = Cl}]
{\small\setlength{\tabcolsep}{3pt}\renewcommand{\arraystretch}{1.25}\begin{center}\begin{tabular}{|L{0.8cm}|L{3.0cm}|L{0.65cm}|L{1.6cm}|}
\hline
\multicolumn{2}{|c|}{\parbox{3.3cm}{\centering List-I\\Dimerization/ polymerization by}} & \multicolumn{2}{c|}{List-II}\\
\hline
(I) & $3c-2e^-$ & (P) & $\mathrm{BeX_2}$\\
\hline
(II) & Co-ordinate bond & (Q) & $\mathrm{AlX_3}$\\
\hline
(III) & Odd $e^-$ & (R) & $\mathrm{CX_3}$\\
\hline
(IV) & No dimerization & (S) & $\mathrm{BX_3}$\\
\hline
\multicolumn{2}{|c|}{} & (T) & $\mathrm{IX_3}$\\
\hline
\multicolumn{2}{|c|}{} & (U) & $\mathrm{ClO_3}$\\
\hline
\end{tabular}\end{center}}
Select correct match when $X=\mathrm{Cl}$
\begin{qoptions}
\opt{A}{I $\to$ P}
\opt{B}{III $\to$ R}
\opt{C}{II $\to$ Q}
\opt{D}{IV $\to$ T}
\end{qoptions}
\end{bankq}
\begin{bankq}[id=TOAS05-P2-CHEM-S2-Q08, source={TOAS - Test 05 (03 Apr 2022)}, paper=2, qno=8, type=numerical, subtopic={Group 13}, answer={0}, answer_src=key, difficulty=3, diff_parts={K1 S1 I1 C0 F1}, diff_note={concepts: structure of the borate anion | steps: identify the sp3 and sp2 boron atoms in the tetraborate ion; atoms forming coordinate bonds | I: structure of the borax anion}]
Total number of boron atom in anionic part of borax which form $\sigma$-coordinate bond is $x$ and total number of boron atom which can form p$\pi$-p$\pi$ coordinate bond is $y$ then calculate of $x-y$
\end{bankq}
\begin{bankq}[id=TOAS05-P2-CHEM-S2-Q09, source={TOAS - Test 05 (03 Apr 2022)}, paper=2, qno=9, type=numerical, subtopic={Group 18}, answer={3}, answer_src=key, difficulty=2, diff_parts={K1 S1 I0 C0 F1}, diff_note={concepts: hydrolysis of oxides and fluorides | steps: identify redox or non-redox hydrolysis for each of the five reactions}]
Find the number of non redox reaction to give their product(s) from the following :\\
(1) $\mathrm{XeF_2+H_2O\to}$\\
(2) $\mathrm{N_2O_5+H_2O\to}$\\
(3) $\mathrm{N_2O_3+H_2O\to}$\\
(4) $\mathrm{XeF_6+H_2O\to}$\\
(5) $\mathrm{XeF_4+H_2O\to}$
\end{bankq}
\begin{bankpassage}[id=TOAS07-P1-CHEM-S3-G5-6, source={TOAS - Test 07 (01 May 2022)}, questions={5-6}]
\bfig{245.1}{TOAS07_p1_chem_s3_q04_1.png}
\end{bankpassage}
\begin{bankq}[id=TOAS07-P1-CHEM-S3-Q05, source={TOAS - Test 07 (01 May 2022)}, paper=1, qno=5, type=numerical, subtopic={Group 15}, answer={1.50}, answer_src=key, difficulty=2, diff_parts={K1 S1 I0 C0 F1}, diff_note={concepts: group 15 compounds from the passage | steps: identify the compounds; ratio X/Y}, group=TOAS07-P1-CHEM-S3-G5-6]
The value of $\left(\dfrac{X}{Y}\right)$ is...........
\end{bankq}
\begin{bankq}[id=TOAS07-P1-CHEM-S3-Q06, source={TOAS - Test 07 (01 May 2022)}, paper=1, qno=6, type=numerical, subtopic={Group 15}, answer={0.25}, answer_src=key, difficulty=2, diff_parts={K1 S1 I0 C0 F1}, diff_note={concepts: group 15 compound bonding | steps: structure of Z; count pi and sigma bonds}, group=TOAS07-P1-CHEM-S3-G5-6]
The ratio of number of $\pi$ bond(s) to number of $\sigma$ bond(s) in `Z' is ...........
\end{bankq}
\begin{bankq}[id=TOAS09-P1-CHEM-S2-Q06, source={TOAS - Test 09 (30 May 2022)}, paper=1, qno=6, type=multi, subtopic={Group 15}, answer={A,B,D}, answer_src=key, difficulty=3, diff_parts={K1 S2 I0 C0 F1}, diff_note={concepts: hydrolysis products of non-metal compounds | steps: P4O6 gives H3PO3; SO2Cl2 gives H2SO4; SF6 inert; SeF6 hydrolyses}]
Which of the following compound(s) give(s) dibasic acid as a product on hydrolysis.
\begin{qoptions}
\opt{A}{P$_4$O$_6$}
\opt{B}{SO$_2$Cl$_2$}
\opt{C}{SF$_6$}
\opt{D}{SeF$_6$}
\end{qoptions}
\end{bankq}
\begin{bankq}[id=TOAS09-P2-CHEM-S2-Q02, source={TOAS - Test 09 (30 May 2022)}, paper=2, qno=2, type=numerical, subtopic={Group 13}, answer={4}, answer_src=solved, working={Dimer Al$_2$(CH$_3$)$_6$ has four terminal Al-C 2c-2e bonds and two 3c-2e bridges}, difficulty=2, diff_parts={K1 S1 I0 C0 F1}, diff_note={concepts: bonding in dimeric alkylaluminium | steps: four terminal Al-C bonds and two bridging 3c-2e bonds}]
The number of 2c -- 2e bonds present in a dimer molecule of Al(CH$_3$)$_3$ is-
\end{bankq}
\end{document}
